\documentclass{article}
\usepackage{graphicx,subcaption}
\title{Figure-Heavy Benchmark: Visual Analysis}
\author{Frank Figures}
\begin{document}
\maketitle

\begin{abstract}
Paper exercising figure resolution, subfigures, cross-references,
and caption parsing.  All asset paths are synthetic and will trigger
\texttt{figure\_missing} warnings when run with \texttt{--no-assets},
which is the intended perf benchmark mode.
\end{abstract}

\section{Architecture Overview}

\begin{figure}[h]
\centering
\includegraphics[width=0.9\textwidth]{figures/architecture_overview.pdf}
\caption{High-level architecture of the proposed system.  The encoder
  processes input tokens $x_1, \ldots, x_n$ and produces hidden states
  $h_1, \ldots, h_n$.}
\label{fig:arch}
\end{figure}

\begin{figure}[h]
\centering
\begin{subfigure}[b]{0.48\textwidth}
\includegraphics[width=\textwidth]{figures/encoder_block.pdf}
\caption{Encoder block detail.}
\label{fig:encoder}
\end{subfigure}
\hfill
\begin{subfigure}[b]{0.48\textwidth}
\includegraphics[width=\textwidth]{figures/decoder_block.pdf}
\caption{Decoder block detail.}
\label{fig:decoder}
\end{subfigure}
\caption{Transformer block internals.}
\label{fig:blocks}
\end{figure}

\section{Attention Visualization}

\begin{figure}[h]
\centering
\includegraphics[width=0.7\textwidth]{figures/attention_heads.png}
\caption{Attention weight matrices for heads 1--4 of layer 6.
  Head 2 attends to syntactic dependencies while head 4 captures
  positional patterns.}
\label{fig:attention}
\end{figure}

\begin{figure}[h]
\centering
\begin{subfigure}[b]{0.32\textwidth}
\includegraphics[width=\textwidth]{figures/attn_layer1.png}
\caption{Layer 1}
\label{fig:attn1}
\end{subfigure}
\hfill
\begin{subfigure}[b]{0.32\textwidth}
\includegraphics[width=\textwidth]{figures/attn_layer6.png}
\caption{Layer 6}
\label{fig:attn6}
\end{subfigure}
\hfill
\begin{subfigure}[b]{0.32\textwidth}
\includegraphics[width=\textwidth]{figures/attn_layer12.png}
\caption{Layer 12}
\label{fig:attn12}
\end{subfigure}
\caption{Attention patterns across layers.}
\label{fig:attn_layers}
\end{figure}

\section{Training Curves}

\begin{figure}[h]
\centering
\includegraphics[width=0.8\textwidth]{figures/loss_curve.pdf}
\caption{Training and validation loss over 100 epochs.  The model
  converges after approximately 60 epochs with early stopping
  patience of 10.}
\label{fig:loss}
\end{figure}

\begin{figure}[h]
\centering
\begin{subfigure}[b]{0.48\textwidth}
\includegraphics[width=\textwidth]{figures/lr_schedule.pdf}
\caption{Learning rate schedule (cosine annealing with warmup).}
\label{fig:lr}
\end{subfigure}
\hfill
\begin{subfigure}[b]{0.48\textwidth}
\includegraphics[width=\textwidth]{figures/gradient_norm.pdf}
\caption{Gradient norm during training.}
\label{fig:grad}
\end{subfigure}
\caption{Training dynamics.}
\label{fig:dynamics}
\end{figure}

\section{Qualitative Results}

\begin{figure}[h]
\centering
\includegraphics[width=\textwidth]{figures/qualitative_grid.png}
\caption{Qualitative comparison of outputs.  Row 1: input images.
  Row 2: ground truth.  Row 3: our method.  Row 4: baseline.}
\label{fig:qual}
\end{figure}

\begin{figure}[h]
\centering
\begin{subfigure}[b]{0.24\textwidth}
\includegraphics[width=\textwidth]{figures/sample_a.png}
\caption{Sample A}
\label{fig:sa}
\end{subfigure}
\hfill
\begin{subfigure}[b]{0.24\textwidth}
\includegraphics[width=\textwidth]{figures/sample_b.png}
\caption{Sample B}
\label{fig:sb}
\end{subfigure}
\hfill
\begin{subfigure}[b]{0.24\textwidth}
\includegraphics[width=\textwidth]{figures/sample_c.png}
\caption{Sample C}
\label{fig:sc}
\end{subfigure}
\hfill
\begin{subfigure}[b]{0.24\textwidth}
\includegraphics[width=\textwidth]{figures/sample_d.png}
\caption{Sample D}
\label{fig:sd}
\end{subfigure}
\caption{Individual sample outputs.}
\label{fig:samples}
\end{figure}

\begin{figure}[h]
\centering
\includegraphics[width=0.6\textwidth]{figures/failure_cases.png}
\caption{Failure cases.  The model struggles with (a) heavily occluded
  objects, (b) rare categories, and (c) extreme lighting conditions.}
\label{fig:failures}
\end{figure}

\section{Cross-Reference Summary}

Our architecture (Figure~\ref{fig:arch}) consists of encoder
(Figure~\ref{fig:encoder}) and decoder (Figure~\ref{fig:decoder}) blocks.
Attention visualizations in Figures~\ref{fig:attention}
and~\ref{fig:attn_layers} reveal interpretable patterns.
Training curves (Figure~\ref{fig:loss}) and dynamics
(Figure~\ref{fig:dynamics}) show stable convergence.
Qualitative results appear in Figures~\ref{fig:qual}
and~\ref{fig:samples}, with failure analysis in Figure~\ref{fig:failures}.

\end{document}
